\documentclass[main]{subfiles}
\usepackage[sols]{workbook}
\numbering{ws2-4}

\begin{document}

\ifSubfilesClassLoaded{%
  \chaptermark{\chaptertwo}
}{}

\section{Definition of Convergence} \label{ws2-4}

\begin{framed}
\centerline{\textbf{Lesson Goals}}

You should be able to:
\begin{itemize}
\item Find the partial sums of a series, and create a partial sum diagram for a series
\item Understand the difference between the \textit{terms} of a series and the \textit{partial sums} of a series, as well as the relationship between them
\item Write an explicit closed formula for the $n$-th partial sum for certain types of series
\item Understand and apply the definition of convergence to determine if a given series converges or diverges
\end{itemize}
\end{framed}




  \begin{enumerate}
    \item 
  \begin{problem}[\vfill]
  Does the series $1 + 1 + 1 + \cdots$ converge or diverge?  If it
  converges, what is its sum?
  \end{problem}

  \begin{solution}
  The $n$-th partial sum $s_n$ is the sum of the first $n$ terms, which
  is just $n$.  Since $\displaystyle \lim_{n \to \infty} s_n = \lim_{n
  \to \infty} n$ does not exist, the series \fbox{diverges}.
  \end{solution}

  \item 
  \note{Sometimes students think the partial sums ``have two limits''
  because half of them tend to 0 and the other half tend to 1.  Make
  sure that they understand that this is still divergent.}
  
  \begin{problem}[\vfill]
  Does the series $1 + (-1) + 1 + (-1) + \cdots$ converge or diverge?
  If it converges, what is its sum?
  \end{problem}

  \begin{solution}
  The partial sums are
  \begin{align*}
  s_1 & = 1 \\
  s_2 & = 1 + (-1) = 0 \\
  s_3 & = 1 + (-1) + 1 = 1 \\
  s_4 & = 1 + (-1) + 1 + (-1) = 0 \\
  & \eqdots
  \end{align*}
  That is, the partial sums alternate between 1 and 0 forever.
  Therefore, the limit of the partial sums, $\displaystyle \lim_{n \to
  \infty} s_n$, does not exist.  This is the same as saying that the
  series \fbox{diverges}.
  \end{solution}
\pspace{\newpage}  
  \begin{framed}
\centerline{\textbf{Partial Sums}}
Let $\displaystyle \sum_{k=1}^{\infty} a_k$ be a series. The \uline{$n$th partial sum} of the series, denoted by $s_n$, is the sum of the first $n$ terms of the series. In particular,
$$\boxed{s_n= a_1 + a_2 + a_3 + \cdots + a_n = \sum_{k=1}^n a_k}$$
Thus, $$\boxed{\displaystyle \sum_{k=1}^{\infty} a_k =\lim_{n \to \infty} \sum_{k=1}^n a_k =  \lim_{n \to \infty} s_n}$$
\vspace{0.5cm}

\centerline{\textbf{The Definition of Convergence of a Series}}
A series $\displaystyle \sum a_k$ \uline{converges} if $\displaystyle \lim_{n \to \infty} s_n$ exists (i.e., approaches a finite value). Otherwise, $\displaystyle \sum a_k$ \uline{diverges}.
\end{framed}


  
  \item 
  \begin{enumerate}
    \item 
    \begin{problem}[2in]
    Does the series $\displaystyle 1 + \frac{1}{2} + \frac{1}{4} +
    \frac{1}{8} + \cdots$ converge or diverge?  If it converges, what is
    its sum?
    \end{problem}

    \begin{solution}
    The $n$-th partial sum $s_n$ is equal to $2 - \frac{1}{2^{n-1}}$;
    one way to see this is by drawing a partial sum diagram like the one
    in (b).  (We will see another way when we study geometric series.)
    The limit $\displaystyle \lim_{n \to \infty} s_n$ of the partial
    sums is equal to 2, which is the same as saying that the series
    \fbox{converges}, and its sum is \fbox{2}.
    \end{solution}

    \item
    \begin{problem}[1in]
    Draw a diagram showing the first 6 partial sums of the series
    $\displaystyle 1 + \frac{1}{2} + \frac{1}{4} + \frac{1}{8} +
    \cdots$.
    \end{problem}

    \begin{solution}
    \begin{center}
    \includegraphics[width=5in]{ws/images/geometric-half-convergence}
    \end{center}
    \end{solution}
  \end{enumerate}

\probonly{\newpage}

  \item 
  \note{Basic practice of using the definition of convergence.}

  In this problem, you'll investigate the series $\displaystyle
  \sum_{k=1}^\infty \frac{2}{(2k)^2 - 1}$.
  \begin{enumerate}
    \item
    \begin{problem}[\vfill]
    Write out the first 4 terms of the series.
    \end{problem}

    \begin{solution}
    $\frac{2}{3}$, $\frac{2}{15}$, $\frac{2}{35}$, $\frac{2}{63}$.
    \end{solution}

    \item
    \begin{problem}[\vfill]
    Find the first 4 partial sums of the series.  Express them as
    simplified fractions.
    \end{problem}

    \begin{solution}
    The first 4 partial sums are $\frac{2}{3}$, $\frac{2}{3} +
    \frac{2}{15}$, $\frac{2}{3} + \frac{2}{15} + \frac{2}{35}$,
    $\frac{2}{3} + \frac{2}{15} + \frac{2}{35} + \frac{2}{63}$.
    Simplified, they are $\frac{2}{3}$, $\frac{4}{5}$, $\frac{6}{7}$,
    $\frac{8}{9}$.  
    \end{solution}

    \item
    \begin{problem}[\vfill]
    Looking at your answer to {{\#4}(b)},
    guess a formula for the $n$-th partial sum.
    \end{problem}

    \begin{solution}
    Based on the pattern of the first 4 partial sums, we guess that the
    $n$-th partial sum is \fbox{$\displaystyle \frac{2n}{2n + 1}$}.
    \end{solution}
    
    \item
    \begin{problem}[\vfill]
    Based on your guess from part (c), does the series converge or
    diverge?  If it converges, what is its sum?
    \end{problem}

    \begin{solution}
    The definition of convergence says we need to look at the limit of
    the partial sums; in this case, that is $\displaystyle \lim_{n \to
    \infty} \frac{2n}{2n + 1}$.  To evaluate this limit, you can use
    L'Hospital's Rule or just divide the top and bottom by the largest
    term ($2n$) to get $\displaystyle \lim_{n \to \infty} \frac{2n}{2n +
    1} = \lim_{n \to \infty} \frac{1}{1 + \frac{1}{2n}} = 1$.  So, the
    series \fbox{converges}, and its sum is \fbox{1}.
    \end{solution}
  \end{enumerate}
  
\pspace{\newpage}
  
  \item 
  \note{This is a series in which it's not so easy (at this point) for
  them to write down a general formula for the partial sums.  So the
  point of using a partial sum diagram is that you can still get a very
  good picture of whether the series converges even without computing
  the partial sums.

  Of course, the disadvantage of this is that you don't know the sum of
  the series, so you may want to talk a bit about why it's useful to
  know whether a series converges even without knowing its sum.  In this
  picture, it's also nice to talk about where they think the sum lies on
  the number line; they should see that it ought to lie between $s_5$
  and $s_6$.
  }
  
    
  \begin{enumerate}
    \item 
    \begin{problem}[\vfill]
    Draw a partial sum diagram showing the first 6 partial sums of the
    series $\displaystyle \sum_{k=0}^\infty \left( - \frac{1}{2}
    \right)^k$.
    \end{problem}

    \begin{solution}
    In $\cdots$ notation, the series is $\displaystyle 1 - \frac{1}{2} +
    \frac{1}{4} - \frac{1}{8} + \cdots$.  Here is the partial sum
    diagram:
    \begin{center}
    \includegraphics[width=5in]{ws/images/partial-sum-diagram-3}
    \end{center}
    \end{solution}

    \item
    \begin{problem}[\vfill]
    Based on your diagram, do you think the series $\displaystyle
    \sum_{k=0}^\infty \left( - \frac{1}{2} \right)^k$ converges or
    diverges?
    \end{problem}

    \begin{solution}
    In drawing the diagram, you might have noticed that each partial sum
    (after the first 2) is halfway between the two previous partial
    sums.  As we keep going, it seems like the partial sums have ``less
    place to go,'' which suggests that the partial sums do indeed have a
    limit (although it's not entirely clear what that limit is).  Since
    the partial sums appear to have a limit, the series should
    \fbox{converge}.
    \end{solution}
    
  \end{enumerate}





  \item 
  \note{The main point here is for them to realize that a term can be
  thought of as the difference between two partial sums.  (This is a
  very useful perspective for problems like {\#6} on the homework.)}

  Here is a diagram showing the first 6 partial sums of an unknown
  series $\displaystyle \sum_{k=1}^\infty a_k$.
  \begin{center}
  \includegraphics[width=5in]{ws/images/info-from-partial-sum-diagram}
  \end{center}

  \begin{enumerate}
    \item
    \begin{problem}[\vfill]
    From the diagram, determine whether each of $a_1, a_2, \ldots, a_6$
    is positive, negative, or zero.
    \end{problem}

    \begin{solution}
    Remember that the $n$-th partial sum is the sum of the first $n$
    terms of the series.  To get from the $(n - 1)$-th partial sum to
    the $n$-th partial sum, we simply add the $n$-th term.  So, the
    $n$-th term is equal to the $n$-th partial sum minus the $(n -
    1)$-th partial sum.  In equations, this tells us that:
    \begin{align*}
    a_2 &= s_2 - s_1 \\
    a_3 &= s_3 - s_2 \\
    a_4 &= s_4 - s_3 \\
    a_5 &= s_5 - s_4 \\
    a_6 &= s_6 - s_5
    \end{align*}
    (Another way to think of these equations is just to write out what
    $s_n$ means: for example, $s_1 = a_1$ and $s_2 = a_1 + a_2$, so it
    is evident that $s_2 - s_1 = a_2$.)

    Looking at the diagram, we see that $s_1$ (which is the same as
    $a_1$) is negative. $a_2 = s_2 - s_1$ is positive (we can see that
    $s_2 > s_1$).  $a_3 = s_3 - s_2$ is positive, $a_4 = s_4 - s_3$ is
    negative, $a_5 = s_5 - s_4$ is positive, and $a_6 = s_6 - s_5$ is
    positive. 
    \end{solution}

    \item
    \begin{problem}[\vfill]
    Which of $a_1, a_2, \ldots, a_6$ is biggest in magnitude?
    \end{problem}

    \begin{solution}
    Using the fact that $a_1 = s_1$, $a_2 = s_2 - s_1$, and so on, the
    biggest in magnitude is $s_4 - s_3 = a_4$. 
    \end{solution}
  \end{enumerate}
  
    \probonly{\newpage}
    
 \item \label{partial-sum-diagram}
\note{Two examples showing that, even if the terms go to 0, the series
need not converge.  These also prepare them to understand why the
harmonic series diverges (similar grouping).}

\begin{enumerate}

  \item \label{partial-sum-diagram-grouping1}
  \begin{problem}
  Draw a diagram showing the first 15 partial
  sums

 of the series $\displaystyle 1 +
  \frac{1}{2} + \frac{1}{2} + \frac{1}{3} + \frac{1}{3} + \frac{1}{3}
  + \underbrace{\frac{1}{4} + \cdots + \frac{1}{4}}_{\textrm{$4$
  terms}} + \underbrace{\frac{1}{5} + \cdots +
  \frac{1}{5}}_{\textrm{$5$ terms}} + \underbrace{\frac{1}{6} +
  \cdots + \frac{1}{6}}_{\textrm{$6$ terms}} + \cdots$.  Does the
  series converge or diverge?  Why?
  \end{problem}
\probonly{\vfill}
  \begin{solution}
  Here is the diagram:
  \begin{center}
  \includegraphics[width=4in]{ws/images/partial-sum-diagram1.pdf}
  \end{center}
  In this case, the series \fbox{diverges} because the partial sums
  tend to infinity.  Why?  Well, let's take a look at some specific
  partial sums.  Suppose we look at the following
  partial sum: 
  \[ \underbrace{1}_{\textrm{1 term}} + \underbrace{\frac{1}{2} +
  \frac{1}{2}}_{\textrm{2 terms}} + \underbrace{\frac{1}{3} +
  \frac{1}{3} + \frac{1}{3}}_{\textrm{3 terms}} + \cdots +
  \underbrace{\frac{1}{100} + \cdots + \frac{1}{100}}_{\textrm{100
  terms}}.\] Each of the bracketed groups sums to 1, and there are
  $100$ groups, so this partial sum is equal to $100$.  How about this
  partial sum?
  \[ \underbrace{1}_{\textrm{1 term}} + \underbrace{\frac{1}{2} +
  \frac{1}{2}}_{\textrm{2 terms}} + \underbrace{\frac{1}{3} +
  \frac{1}{3} + \frac{1}{3}}_{\textrm{3 terms}} + \cdots +
  \underbrace{\frac{1}{200} + \cdots + \frac{1}{200}}_{\textrm{200
  terms}}\]
  Again, each bracketed group sums to 1, but now there are 200 groups,
  so this partial sum is equal to 200.  Similarly, the partial sum
  \[ \underbrace{1}_{\textrm{1 term}} + \underbrace{\frac{1}{2} +
  \frac{1}{2}}_{\textrm{2 terms}} + \underbrace{\frac{1}{3} +
  \frac{1}{3} + \frac{1}{3}}_{\textrm{3 terms}} + \cdots +
  \underbrace{\frac{1}{300} + \cdots + \frac{1}{300}}_{\textrm{300
  terms}}\]
  is equal to 300.  As you can see, we can continue in this way to
  make partial sums as big as we like, which shows that the partial
  sums have no limit.
  \end{solution}

  \begin{grading}
  They should say that the partial sums tend to $\infty$, but they
  don't need to give a detailed explanation of why that is.
  \end{grading}

  \item \label{partial-sum-diagram-grouping2}
  \begin{problem}
  Does the series $\displaystyle 1 - \frac{1}{2} - \frac{1}{2} +
  \frac{1}{3} + \frac{1}{3} + \frac{1}{3} - \underbrace{\frac{1}{4} -
  \cdots - \frac{1}{4}}_{\textrm{4 terms}} + \underbrace{\frac{1}{5}
  + \cdots + \frac{1}{5}}_{\textrm{5 terms}} -
  \underbrace{\frac{1}{6} - \cdots - \frac{1}{6}}_{\textrm{6 terms}} +
  \cdots$ converge or diverge?  Why?  (It will be helpful to draw a
  diagram like in part (a), but you need not turn it in since it's
  hard to make the diagram neat.)
  \end{problem}
\probonly{\vfill}

  \begin{solution}
  The first partial sum is just 1.  The next 2 partial sums are:
  \begin{center}
  \includegraphics[width=4in]{ws/images/partial-sum-diagram2.pdf}
  \end{center}
  That is, they decrease to 0.  The next 3 partial sums after that
  are: 
  \begin{center}
  \includegraphics[width=4in]{ws/images/partial-sum-diagram3.pdf}
  \end{center}
  That is, they increase back to 1.  The next 4 partial sums are:
  \begin{center}
  \includegraphics[width=4in]{ws/images/partial-sum-diagram4.pdf}
  \end{center}
  They've decreased back to 0.  The next 5 partial sums are:
  \begin{center}
  \includegraphics[width=4in]{ws/images/partial-sum-diagram5.pdf}
  \end{center}
  They've increased back to 1 again.  This pattern continues, the
  partial sums increasing to 1 and decreasing back to 0 forever.
  Therefore, the partial sums have no limit, and the series
  \fbox{diverges}. 
  \end{solution}

\end{enumerate}

  \end{enumerate}
  
\end{document}